\chapter{Sensores}
La medición de la calidad del aire se apoya en la evaluación y cuantificación de la presencia y concentración de contaminantes y componentes atmosféricos que pueden afectar la salud humana, el medio ambiente y los ecosistemas. Esta medición se realiza utilizando una variedad de sensores especializados para monitorear diferentes parámetros y contaminantes presentes en el aire. La medición de la calidad del aire se realiza a través de redes de monitoreo, que consisten en estaciones fijas distribuidas estratégicamente en diferentes ubicaciones geográficas. También se utilizan sensores portátiles y dispositivos de monitoreo en tiempo real para obtener mediciones más detalladas y específicas en áreas particulares o durante estudios de campo.

El análisis de los datos recolectados permite evaluar la calidad del aire, identificar fuentes de contaminación, evaluar la eficacia de las medidas de control y tomar decisiones informadas para proteger la salud y el medio ambiente. Además, estos datos son fundamentales para establecer políticas y regulaciones ambientales que promuevan una mejor calidad del aire y un entorno más saludable.

\section{Clasificación de sensores para calidad del aire}

\textls[-10]{Entre las principales variables que se evalúan dentro de lo que se considera calidad del aire, se encuentran los contaminantes gaseosos, las partículas suspendidas, los contaminantes biológicos y los parámetros meteorológicos.}

\subsection{Contaminantes gaseosos}

Dentro de esta denominación se incluyen gases como el dióxido de carbono (CO\textsubscript{2}), óxidos de nitrógeno (NO\textsubscript{2}), el ozono (O\textsubscript{3}), el dióxido de azufre (SO\textsubscript{2}) y algunos compuestos orgánicos volátiles (COV), entre otros. Estos contaminantes se miden para determinar su concentración y evaluar si están dentro de los límites aceptables establecidos por las normativas ambientales. Un esquema con los elementos y los principios básicos de su medición se muestra en la figura \ref{fig:fig1A}.

\begin{figure}[ht]
  \begin{center}
    \caption{Esquema básico de un sensor para contaminantes gaseosos \cite{sengascon}}
    \label{fig:fig1A}
    \includegraphics[width=.8\textwidth]{Fig1A}
    \propia
  \end{center}
\end{figure}

Los principales sensores para contaminantes gaseosos incluyen:

\begin{itemize}
  \item Sensor de dióxido de carbono (\CO): mide la concentración de dióxido de carbono en el aire. Es utilizado para controlar la ventilación en edificios, monitorear la calidad del aire en interiores y detectar posibles fuentes de contaminación.
  \item Sensor de COV: mide la presencia de compuestos químicos volátiles en el aire, como los emitidos por productos químicos, solventes, productos de limpieza y materiales de construcción. Los COV pueden contribuir a la mala calidad del aire y afectar la salud humana.
  \item Sensor de óxidos de nitrógeno (\NOX): mide la concentración de óxidos de nitrógeno en el aire, incluyendo el dióxido de nitrógeno ($\textrm{NO}_2$) y el monóxido de nitrógeno ($\textrm{NO}$). Estos gases son emitidos principalmente por la combustión de combustibles fósiles y pueden tener efectos negativos en la calidad del aire y la salud.
  \item Sensor de ozono ($\textrm{O}_3$): mide la concentración de ozono en el aire. El ozono puede ser perjudicial para la salud humana en altas concentraciones y es un contaminante común en áreas urbanas y cerca de fuentes de emisión.
\end{itemize}

\subsection{Partículas suspendidas}

Corresponde a la concentración de partículas, tanto finas (PM2.5) como grandes (PM10), en el aire. Estas partículas pueden provenir de fuentes naturales o antropogénicas, como la quema de combustibles fósiles, incendios forestales, el polvo de construcción, entre otros. Dentro de la familia de sensores para la medición de este tipo de afectación del aire, se tienen los siguientes:

\begin{itemize}
  \item Sensor de dispersión láser: utiliza un haz de luz láser para medir el tamaño y la concentración de partículas suspendidas. Mide la dispersión de la luz causada por las partículas en el aire y proporciona información sobre el tamaño de estas y su distribución.
  \item Sensor de impacto: este tipo de sensor utiliza una placa o sustrato sensible al impacto de las partículas. Cuando las partículas chocan contra la superficie del sensor, se genera una señal eléctrica proporcional a la concentración de partículas.
  \item Sensor de efecto túnel: este tipo de sensor se basa en el principio del efecto túnel cuántico, en el que las partículas suspendidas en el aire pueden alterar la corriente que fluye a través de una barrera de túnel. La alteración de la corriente se detecta y se utiliza para medir la concentración de estas.
  \item Sensor de masa: este sensor utiliza un filtro o un medio poroso para atrapar las partículas suspendidas en el aire. Luego, la masa de las partículas capturadas se pesa y se determina la concentración de estas en función de la masa medida.
  \item Sensor de condensación de vapor: este sensor utiliza la técnica de condensación de vapor para aumentar el tamaño de las partículas y hacerlas detectables. El vapor se condensa sobre las partículas, haciéndolas crecer en tamaño, y luego se mide el tamaño de las partículas condensadas.
\end{itemize}

\subsection{Contaminantes biológicos}

Se pueden medir microorganismos como bacterias, hongos, polen y esporas, entre otros microorganismos, que pueden representar riesgos para la salud humana. Algunos de los tipos de sensores utilizados para la detección de contaminantes biológicos son:

\begin{itemize}
  \item Sensores de PCR (reacción en cadena de la polimerasa): utilizan la técnica de PCR para amplificar y detectar material genético específico de los organismos objetivo. Estos sensores son altamente sensibles y se utilizan comúnmente para la detección de virus, bacterias y otros patógenos.
  \item \textls[-10]{Sensores de inmunodetección: utilizan anticuerpos u otras mo\-lécu\-las de reconocimiento específicas para detectar la presencia de agentes biológicos. Estos sensores pueden detectar proteínas o componentes celulares específicos y se utilizan en aplicaciones como la detección de toxinas bacterianas o virus.}
  \item Sensores de microarrays de ADN: utilizan una matriz de sondas de ADN específicas para detectar y distinguir diferentes secuencias de ADN. Estos sensores son útiles para la identificación de microorganismos y pueden detectar múltiples especies o cepas en una sola prueba.
  \item Sensores de espectroscopia: utilizan técnicas de espectroscopia para detectar y analizar la firma espectral de los contaminantes biológicos. Estos sensores pueden utilizar diferentes rangos del espectro electromagnético, como la luz visible, infrarroja o ultravioleta, para identificar características específicas de los organismos o sus productos de desecho.
  \item Sensores basados en biosensores: utilizan sistemas biológicos, como enzimas o células vivas, para detectar la presencia de contaminantes biológicos. Estos sensores pueden detectar cambios en la actividad enzimática o respuestas biológicas específicas para identificar la presencia de microorganismos o sustancias biológicas.
\end{itemize}

\subsection{Parámetros meteorológicos}

Además de los contaminantes específicos, también se miden variables meteorológicas como la temperatura, la humedad relativa y la velocidad y dirección del viento. Estos datos son importantes para comprender los patrones de dispersión de los contaminantes y su impacto en diferentes áreas. Algunos sensores destacados para estas magnitudes son:

\begin{itemize}
  \item Sensor de temperatura: mide la temperatura ambiente y puede utilizar termistores, termopares, sensores de resistencia (RTD), sensores semiconductores y sensores metálicos.
  \item Sensor de humedad relativa: mide la cantidad de humedad en el aire, de forma indirecta, por medición directa de la capacitancia, la resistencia o la conductancia eléctrica del medio húmedo.
  \item Sensor de presión atmosférica: mide la presión atmosférica, indirectamente, utilizando medios actinométricos sobre cuerpos de prueba como membranas, fuelles, tubos, entre otros.
  \item Sensores de lluvia o pluviómetros: miden la cantidad de precipitación que cae en un determinado periodo de tiempo. Los pluviómetros generalmente consisten en un recipiente graduado que recoge y mide el agua de lluvia, en un área y un tiempo determinados, para dar base a la proyección de cálculos asociados al volumen y caudal totales obtenibles de dichos parámetros.
\end{itemize}

\section[Sensores específicos para contaminantes gaseosos]{Sensores específicos para\\ contaminantes gaseosos}

\subsection{Sensores para CO\textsubscript{2}}

Las técnicas de detección de \CO\ han recorrido un largo camino en el desarrollo de sensores para el análisis de gases, pasando por métodos químicos y físicos como la cromatografía, el análisis infrarrojo, la espectrometría, la conductimetría, entre otras \cite{Kumar-2007}.

Actualmente se manejan, fundamentalmente, dos tipos de técnicas para la detección de \CO:
\begin{itemize}
  \item Determinación de la concentración de CO\textsubscript{2} en el aire por el método \gls{ndir}, el cual es el más empleado para la detección en tiempo real.
  \item Detección por medio de sensores electrolíticos de estado sólido.
\end{itemize}

\subsubsection{Sensor infrarrojo no dispersivo}
Un sensor \gls{ndir} es un dispositivo espectroscópico simple de uso frecuente como detector de gas.
Los componentes principales de este sensor se presentan en la figura \ref{fig:fig2A} y son los siguientes \cite{Yi-2005}:
\begin{itemize}
  \item Una lámpara de IR que se constituye en la fuente \gls{ndir}. Es el sistema \gls{ndir} propiamente dicho.
  \item Una cámara o tubo de luz por el cual circula la muestra.
  \item Un filtro de IR.
  \item Un detector de luz IR.
\end{itemize}

El gas se difunde en la cámara y la concentración de \CO\ se mide ópticamente por la absorción de una longitud de onda específica en el rango del infrarrojo. La luz IR se dirige, atravesando la cámara que contiene la muestra, hacia el detector; un filtro óptico, en frente del detector, elimina todas las componentes (longitudes de onda) de la luz, excepto la longitud de onda que absorben las moléculas del gas seleccionado. El \CO\ presenta una fuerte absorbancia en la banda IR; la longitud de onda frecuentemente empleada es de \SI{4260}{nm}.

\begin{figure}[H]
  \begin{center}
    \caption{Sensor \gls{ndir} para \CO}
    \label{fig:fig2A}
    \includegraphics[width=0.9\textwidth]{Fig2A}
    \propia
  \end{center}
\end{figure}

El término no dispersivo se refiere a que toda la luz (todas las longitudes de onda) atraviesa la muestra de gas y el filtrado se aplica solamente a la luz que logra llegar al final de la cámara, que es cuando esta ingresa al detector.

\subsubsection{Función característica}

La absorción de la radiación IR por parte de las moléculas de \CO\ sigue la ley de Beer-Lambert, que se expresa así:
\begin{align}
  A&=\log_{10}(I_a/I_f)\\
  A&=a \times x \times C
  \label{Lamb}
\end{align}
Donde:

\begin{itemize}
  \item[$A$] $\rightleftharpoons$ absorbancia del \CO. Es una magnitud adimensional.
  \item[$I_a$] $\rightleftharpoons$ Intensidad de la luz resultante en cualquier punto del sistema a lo largo de la cámara. Se mide en W/sr.
  \item[$I_f$] $\rightleftharpoons$ Intensidad de la luz que emerge de la fuente. Se mide en W/sr.
  \item[$a$] $\rightleftharpoons$ Coeficiente de absorción del \CO; se denomina igualmente como \emph{absortividad molar}. Se mide en L/(mol$\cdot$cm).
  \item[$x$] $\rightleftharpoons$ Longitud de la trayectoria, camino óptico recorrido o longitud de la celda. Se mide en cm.
  \item[$C$] $\rightleftharpoons$ Concentración del \CO. Se mide en mol/L.
\end{itemize}

La intensidad de luz, $I$, que llega al sensor óptico, disminuye en la medida en que se aumenta la concentración, $C$, de \CO\ en la cámara. Si la longitud de paso óptico de la cámara es $L$ y la absortividad molar del \CO\ es $A$, la intensidad de luz en el sensor óptico queda definida por la ley de Beer-Lambert:
\begin{equation*}
  I=I_0\times 10^{CLA}
\end{equation*}
Donde $I_0$ es la luz a la salida de la fuente.

\subsection{Sensores para COV}
Un sensor de COV es ampliamente utilizado en aplicaciones de monitoreo ambiental, control de calidad del aire y detección de gases en cualquier entorno.
El funcionamiento de un sensor de COV se basa en la interacción de los compuestos orgánicos volátiles con una capa sensible en el sensor. Cuando los COV están presentes en el aire, interactúan con la capa sensible, lo que produce cambios en sus propiedades eléctricas, ópticas o químicas. Estos cambios son medidos por el sensor y se convierten en una señal eléctrica que puede ser interpretada como la concentración de COV presente en el ambiente.

El modelo matemático de un sensor de COV varía según el tipo de sensor y la tecnología empleada. Algunos sensores de COV utilizan resistencias eléctricas, conductancias iónicas o espectroscopia para medir los cambios en la denominada capa sensible, y así poder calcular la concentración de COV.
Normalmente, cada fabricante de sensores para COV puede tener su propio modelo matemático, basado en el diseño y la tecnología utilizados en su sensor específico. Los modelos más detallados pueden involucrar ecuaciones racionales o empíricas, basadas en parámetros o en la respuesta del sistema, respectivamente.
El sensor de COV consta de varios componentes principales que trabajan juntos para detectar y medir la concentración de COV en el aire. Estos componentes son los siguientes:
\begin{itemize}
  \item La capa sensible: es la parte del sensor que entra en contacto con los COV presentes; está diseñada para interactuar con los COV y experimentar cambios en sus propiedades físicas o químicas cuando se expone a ellos.
  \item El transductor: es el elemento que convierte los cambios en las propiedades de la capa sensible en una señal eléctrica medible. Dependiendo del tipo de sensor, el transductor puede ser un resistor, un condensador, un diodo o algún otro dispositivo que responda a los cambios en la capa sensible.
  \item Un circuito de acondicionamiento de señal: es el circuito electrónico que amplifica y procesa la señal eléctrica generada por el transductor. Este circuito puede incluir amplificadores, filtros y otros componentes para garantizar una lectura precisa y estable de la concentración de COV.
  \item La fuente de alimentación: proporciona la energía necesaria para el funcionamiento de todo el sistema.
  \item La interfaz de salida: es la parte del sensor que permite la conexión con otros dispositivos o sistemas para transmitir la información sobre la concentración de COV. Esto puede incluir interfaces analógicas, como voltaje o corriente, o interfaces digitales, como un bus de comunicación serial.
\end{itemize}
Es importante destacar que los componentes específicos y la configuración del sensor pueden variar según el diseño y la tecnología utilizada. Cada fabricante puede tener su propia implementación y selección de componentes para lograr el mejor rendimiento y precisión en la medición de COV.

\subsection{Sensores electroquímicos para \texorpdfstring{NO\boldmath{$_{x}$}}{NOx}}

Los sensores electroquímicos para óxido de nitrógeno (\NOX) se basan en reacciones electroquímicas que ocurren en una celda sensora. Estos sensores constan de tres componentes principales: un electrodo de trabajo, un electrodo de referencia y un contraelectrodo. A continuación se presenta un modelo matemático híbrido (racional-empírico) utilizado en este tipo de sensores.

La ley de Nernst establece la relación entre el potencial eléctrico y la concentración de especies químicas en una celda electroquímica. El modelo matemático para un sensor de \NOX\ puede ser expresado así:
%%%%inicioecuacion
\begin{gather}
  \label{nox}
  E=E_o-(RT/nF)\times \ln(C)
\end{gather}
Donde:

\begin{itemize}
  \item[$E$] $\rightleftharpoons$  Potencial eléctrico medido en el electrodo.
  \item[$E_o$] $\rightleftharpoons$  Potencial de base de la reacción electroquímica.
  \item[$R$] $\rightleftharpoons$  Constante universal de los gases ideales.
  \item[$T$] $\rightleftharpoons$  Temperatura termodinámica.
  \item[$F$] $\rightleftharpoons$  Constante de Faraday.
  \item[$C$] $\rightleftharpoons$  Concentración del \NOX.
\end{itemize}

Este modelo establece una relación logarítmica entre el potencial eléctrico medido y la concentración de \NOX.
Existe, también, un modelo denominado de difusión y reacción; en dicho modelo se considera la difusión de los gases de \NOX\ desde el entorno hacia la superficie del electrodo, donde ocurren las reacciones electroquímicas. El modelo matemático puede incluir las ecuaciones de difusión de Fick y reacciones de oxidación y reducción en el electrodo para describir el comportamiento del sensor.
Otros modelos más complejos pueden considerar parámetros adicionales, como el coeficiente de difusión, la velocidad de transferencia de electrones y la resistencia interna del sensor, para obtener una descripción más precisa del comportamiento del sensor electroquímico para \NOX.
Los modelos matemáticos específicos pueden variar según el diseño y la tecnología del sensor electroquímico utilizado. Los fabricantes de sensores electroquímicos suelen proporcionar documentación técnica que incluye información sobre el modelo matemático y las ecuaciones asociadas utilizadas en sus sensores específicos.

\section{Sensores específicos para material particulado}

\subsection{Sensor de dispersión láser para material particulado}

El sensor de dispersión láser para material particulado mantiene la técnica de dispersión de un haz de luz que atraviesa el material por analizar, sufriendo diferentes niveles de atenuación en sus componentes de frecuencia, debido a las diferentes partículas encontradas en su camino. Los elementos que lo componen y sus funciones individuales se pueden resumir así:
\begin{itemize}
  \item Fuente láser: puede ser un láser, generalmente de infrarrojo, con un espectro amplio.
  \item Cámara de dispersión: permite el encuentro entre el láser y el aire contaminado con las partículas en cuestión.
  \item Detector de partículas: los métodos comunes incluyen la dispersión de luz debido al encuentro entre partículas y el haz de luz.
  \item Contador de partículas: una vez que las partículas son detectadas, el sensor captura y cuenta las partículas presentes en una muestra de aire. Esto se puede lograr mediante técnicas como la atracción electrostática, la filtración o la interacción de las partículas con una superficie sensible.
  \item Conversor de partículas a señal eléctrica: el sensor convierte la información de las partículas detectadas en una señal eléctrica, generalmente corriente.
\end{itemize}
Un sensor típico de esta tecnología es el producido por Honeywell, cuyo esquema simplificado se presenta en la figura \ref{fig:fig3A}.
Es importante tener en cuenta que los sensores de material particulado pueden utilizar diferentes tecnologías y principios de funcionamiento, como la dispersión láser, la captación en filtros o la medición de la conductividad.

Un posible modelo matemático para un sensor de partículas de dispersión de luz es el siguiente:
\begin{align}
  C = k \times (I / D)
\end{align}

Donde:
\begin{itemize}
  \item[$C$] $\rightleftharpoons$ Concentración de partículas en el aire.
  \item[$I$] $\rightleftharpoons$ Intensidad de la luz dispersada o atenuada por las partículas.
  \item[$D$] $\rightleftharpoons$ Distancia entre la fuente de luz y el detector.
  \item[$k$] $\rightleftharpoons$ Sensibilidad de la cámara de dispersión.
\end{itemize}

Este modelo asume una relación lineal entre la intensidad de la luz dispersada y la concentración de partículas. La sensibilidad de la cámara se determina experimentalmente y puede variar dependiendo de las características específicas del sensor y del tipo de partículas que se estén midiendo.

\begin{figure}[H]
  \centering
  \caption{Sensor de dispersión para material particulado}
  \label{fig:fig3A}
  \includegraphics[width=0.7\textwidth]{Fig3A}

  \propia
\end{figure}

\subsubsection[Sensor de dispersión láser multiángulo para material particulado]{Sensor de dispersión láser multiángulo\\ para material particulado}

Un novedoso sensor de dispersión para partículas PM2.5 es el desarrollado en el State Key Laboratory of Industrial Control Technology, del College of Control Science and Engineering, de la Universidad de Zhejiang, el cual permite eliminar los efectos de factores meteorológicos, como la humedad, en las partículas. Se concentra en detectar la dispersión en tres ángulos diferentes: \SI{40}{\degree}, \SI{55}{\degree} y \SI{140}{\degree}. Un esquema con arquitectura similar a la disposición real del sensor se presenta en la figura \ref{fig:fig4A}.

\begin{figure}[h]
  \centering
  \caption{Sensor de dispersión multiángulo para material particulado}
  \label{fig:fig4A}
  \includegraphics[width=\textwidth]{Fig4A}

  \propia
\end{figure}

\section{Interfaces de sensores}

\subsection{Sensores con interfaz análogica}

Cuando un sensor presenta un cambio en sus propiedades eléctricas proporcional a los cambios en su variable de medida, es recomendable implementar una cadena de medida que permita su correcto funcionamiento y a su vez un aprovechamiento máximo del rango completo de detección del \gls{adc} del microcontrolador. También deberá considerarse si es necesario hacer algún tipo de tratamiento básico a los datos (promedio, filtrado, entre otros) antes de concluir un valor que se pueda enviar por la red LoRaWAN.

\subsection{Sensores con interfaz digital}

La estructura general consiste en el envío de uno o más comandos, con el fin de obtener una respuesta a estos. El formato de cada comando y su respuesta asociada depende fuertemente del fabricante de cada sensor en particular.
